\subsection{原子、核、粒子与凝聚态补充}

\begin{结论}{Moseley 定律与 X 射线谱}
K$\alpha$ 线（$n=2\to1$，内层电子屏蔽一个核电荷）：
\[
\重点{E_{K\alpha}\approx\frac34\times13.6\text{ eV}\times(Z-1)^2}=10.2\text{ eV}\,(Z-1)^2 .
\]
Cu：$8.0$ keV（实测 $8.05$）。X 光管连续谱的短波限（Duane--Hunt）$\lambda_{\min}=\frac{hc}{eV}=\frac{1.24\text{ nm}}{V[\text{kV}]}$，可测 $\frac he$。荧光产额随 $Z^4$ 增大，轻元素以 Auger 电子为主。
\end{结论}

\begin{结论}{Lamb 位移与电子反常磁矩}
Dirac 理论中 $2S_{1/2}$ 与 $2P_{1/2}$ 简并；QED 真空涨落使 $2S_{1/2}$ 上移约 \重点{1057.8\text{ MHz}}（Lamb--Retherford 1947）。
电子 $g=2(1+a_e)$，$a_e=\frac{\alpha}{2\pi}+\cdots\approx0.00115965$，理论与实验符合到 $10^{-12}$ 量级（物理学最精确的检验之一）。$\mu$ 子 $a_\mu$ 对重粒子圈图更敏感（$\propto m^2$）。
\end{结论}

\begin{结论}{Franck--Condon 原理与振动光谱}
电子跃迁极快，核位置不变（竖直跃迁），振动带强度 $\propto|\langle v'|v''\rangle|^2$（Franck--Condon 因子）；激发态平衡核间距变化越大，振动进程越长。
典型振动波数：H$_2$ $4401$ cm$^{-1}$，HCl $2991$，CO $2143$，O$_2$ $1580$ cm$^{-1}$；$\tilde\nu=\frac1{2\pi c}\sqrt{\frac k\mu}$（同位素 \重点{\propto\mu^{-1/2}}，D 替换 H 约降为 $\frac1{\sqrt2}$）。
\end{结论}

\begin{结论}{Raman 散射：Stokes 与反 Stokes}
激发光 $\nu_0$，振动 $\nu_v$：Stokes $\nu_0-\nu_v$、反 Stokes $\nu_0+\nu_v$。强度比
\[
\重点{\frac{I_{AS}}{I_S}=\Bigl(\frac{\nu_0+\nu_v}{\nu_0-\nu_v}\Bigr)^4\ee^{-h\nu_v/k_BT}}
\]
（Boltzmann 布居 $\times$ $\nu^4$ 散射律），可作非接触测温。Raman 活性要求极化率随振动变化，红外活性要求偶极矩变化；有对称中心的分子二者互斥。
\end{结论}

\begin{结论}{原子核的尺度与结合能}
$R\approx1.2A^{1/3}$ fm，核密度 $\approx2.3\times10^{17}$ kg/m$^3$（与 $A$ 无关，饱和）。比结合能约 $8$ MeV，峰在 $^{62}$Ni/$^{56}$Fe（$8.79$ MeV）。
幻数 \重点{2,8,20,28,50,82,126}：需要强自旋--轨道耦合（Mayer--Jensen 壳模型）；双幻核 $^4$He、$^{16}$O、$^{40}$Ca、$^{208}$Pb 特别稳定。核子 Fermi 动量 $\approx250$ MeV/$c$。
\end{结论}

\begin{结论}{$\alpha$ 衰变：Geiger--Nuttall 律}
Gamow 隧穿因子 $\ee^{-2G}$，$G\approx\frac{\pi Z_1Z_2e^2}{4\pi\varepsilon_0\hbar v}$（远低于势垒时）：
\[
\重点{\log_{10}t_{1/2}=a\frac{Z}{\sqrt Q}+b}.
\]
$Q$ 从 $4$ MeV 增到 $9$ MeV，半衰期从 $10^{17}$ s 降到 $10^{-6}$ s——跨 $23$ 个量级。同样的隧穿因子控制恒星中的聚变（Gamow 峰）。
\end{结论}

\begin{结论}{放射性测年}
母核 $P$、子核 $D$（初始无子核）：
\[
\重点{t=\frac1\lambda\ln\Bigl(1+\frac DP\Bigr)}=\frac{t_{1/2}}{\ln2}\ln\Bigl(1+\frac DP\Bigr).
\]
$^{14}$C：$t_{1/2}=5730$ 年，适用 $\lesssim5\times10^4$ 年；U--Pb、K--Ar（$1.25$ Gyr）、Rb--Sr 等时线法（多个样品回归，不需初始子核为零）。等时线：$\frac{D}{D_s}=\frac{D_0}{D_s}+\frac P{D_s}(\ee^{\lambda t}-1)$，斜率给年龄。
\end{结论}

\begin{结论}{Breit--Wigner 共振截面}
\[
\重点{\sigma(E)=\frac{2J+1}{(2s_1+1)(2s_2+1)}\frac{\pi}{k^2}\frac{\Gamma_i\Gamma_f}{(E-E_0)^2+\Gamma^2/4}},\qquad \Gamma=\frac\hbar\tau .
\]
峰值受幺正性限制 $\sim\frac{4\pi}{k^2}(2l+1)$。$Z$ 玻色子：$M=91.19$ GeV，$\Gamma=2.495$ GeV，由宽度数出轻中微子代数 $N_\nu=2.98\pm0.01$。
\end{结论}

\begin{结论}{$\mu$ 子衰变与 Fermi 常数}
\[
\重点{\Gamma_\mu=\frac{G_F^2m_\mu^5}{192\pi^3}}\ (\hbar=c=1),\qquad \tau_\mu=2.197\ \mu\text{s}\ \Rightarrow\ G_F=1.166\times10^{-5}\text{ GeV}^{-2}.
\]
$\tau$ 轻子同一公式乘衰变道数（$e,\mu$ 与三色夸克 $\approx5$）：$\tau_\tau\approx\tau_\mu\bigl(\frac{m_\mu}{m_\tau}\bigr)^5\times\frac15\approx3\times10^{-13}$ s（实测 $2.9\times10^{-13}$ s）。弱相互作用“弱”是因为 $G_FE^2\ll1$。
\end{结论}

\begin{结论}{截面、亮度与事例率}
事例率 $\dot N=\sigma\mathcal L$；$1$ b $=10^{-28}$ m$^2$，$1$ fb$^{-1}$ 积分亮度对 $1$ fb 截面给 $1$ 个事例。
LHC 设计亮度 $10^{34}$ cm$^{-2}$s$^{-1}$：总非弹截面 $\approx80$ mb，每秒 $\sim10^9$ 次碰撞；Higgs 产生截面 $\sim50$ pb（$13$ TeV），每 fb$^{-1}$ 约 $5\times10^4$ 个。固定靶：$\mathcal L=$ 束流粒子率 $\times$ 靶面密度。
\end{结论}

\begin{结论}{BCS 超导的普适比}
弱耦合 BCS：\重点{2\Delta(0)=3.53k_BT_c}；比热跃变 $\frac{\Delta C}{\gamma T_c}=1.43$；同位素效应 $T_c\propto M^{-1/2}$（声子机制的证据）。
$T_c$ 附近 $\Delta(T)\approx3.06k_BT_c\sqrt{1-T/T_c}$。强耦合材料（Pb、Hg）比值偏大；铜氧化物高温超导 $d$ 波，比值约 $5$--$8$。
\end{结论}

\begin{结论}{电学量子标准：Josephson 与量子 Hall}
Josephson 频率 \重点{f=\frac{2e}hV}\ $=483.6$ GHz/mV；Shapiro 台阶 $V_n=n\frac{hf}{2e}$。
量子 Hall：$\sigma_{xy}=\nu\frac{e^2}h$，von Klitzing 常数 $R_K=\frac h{e^2}=25812.807\ \Omega$。二者使电压、电阻直接溯源到 $e,h$（2019 SI 中二者为定义值）；结合电子泵/电流测量构成“量子计量三角”。
\end{结论}

\begin{结论}{铁磁自旋波与 Bloch $T^{3/2}$ 律}
Heisenberg 铁磁体低能激发为磁振子，色散 \重点{\hbar\omega=Dk^2}（二次，与声子的线性不同）。热激发磁振子数 $\propto T^{3/2}$：
\[
\frac{M_s(0)-M_s(T)}{M_s(0)}\propto T^{3/2},\qquad C_{\text{mag}}\propto T^{3/2}.
\]
反铁磁磁振子线性色散，$C\propto T^3$。一维、二维各向同性 Heisenberg 模型有限温度无长程序（Mermin--Wagner）。
\end{结论}

\begin{结论}{电子偶素与 $\mu$ 子原子}
电子偶素：约化质量 $\frac{m_e}2$，结合能 $6.8$ eV，基态单态（para）衰变为 $2\gamma$，寿命 \重点{125\text{ ps}}；三重态（ortho）衰变为 $3\gamma$，寿命 $142$ ns。
$\mu$ 子氢：玻尔半径缩小 $\approx186$ 倍（约化质量），对质子大小敏感 $\sim186^3$ 倍：$\mu$H Lamb 位移测得的质子半径 $0.84$ fm（“质子半径之谜”，后来电子测量也趋向此值）。
\end{结论}

\begin{结论}{核磁矩与 NMR 频率}
$\omega=\gamma B$：质子 $\frac\gamma{2\pi}=42.58$ MHz/T，$^{13}$C $10.71$，$^{19}$F $40.05$，$^{31}$P $17.24$ MHz/T；电子 $28.0$ GHz/T。
热平衡极化 $P\approx\frac{\hbar\gamma B}{2k_BT}$：$1.5$ T、室温质子约 $5\times10^{-6}$——MRI 信号弱的根源（靠 $10^{23}$ 个自旋弥补）。化学位移以 ppm 计（$\sim10^{-6}$），需 $B$ 均匀到 $10^{-9}$。
\end{结论}

\begin{结论}{固体中的典型能量尺度}
Fermi 能 $\sim5$ eV（$T_F\sim5\times10^4$ K），Debye 能 $\sim20$ meV（$\Theta_D\sim300$ K），带隙 Si $1.12$ eV、GaAs $1.42$ eV、金刚石 $5.5$ eV，激子结合 $\sim10$ meV（GaAs $4.2$ meV），超导能隙 $\sim1$ meV，磁有序 $10$--$1000$ K。
\重点{\text{室温 }k_BT=25.9\text{ meV}}：与这些尺度比较即可判断哪些自由度被激活。光子能量 $E[\text{eV}]=\frac{1240}{\lambda[\text{nm}]}$。
\end{结论}

\begin{结论}{Wiedemann--Franz 之外：Fermi 液体与电子比热}
金属电子比热 $C=\gamma T$，$\gamma=\frac{\pi^2}3k_B^2g(E_F)$：只有 $\sim\frac{k_BT}{E_F}$ 的电子参与。重费米子化合物 $\gamma$ 大 $100$--$1000$ 倍（有效质量大）。
低温 $C=\gamma T+\beta T^3$：作 $\frac CT$ 对 $T^2$ 图，截距 $\gamma$、斜率 $\beta$（得 $\Theta_D$）——\重点{\text{分离电子与声子贡献的标准方法}}。
\end{结论}

\begin{结论}{中微子与粒子物理数量级}
太阳中微子通量 $6.5\times10^{10}$ cm$^{-2}$s$^{-1}$（地表）；MeV 中微子与物质截面 $\sim10^{-44}$ cm$^2$，铅中平均自由程约 $10$ 光年量级。
$E\ll M_W$ 时 $\sigma\sim G_F^2s$，固定靶 $s\propto E_\nu$，所以\重点{\text{截面随实验室能量线性增长}}：$\sigma_{\nu N}\approx0.7\times10^{-38}\,\text{cm}^2\times\frac{E_\nu}{\text{GeV}}$（带电流，每核子）。PeV 能区才使地球对中微子不透明——高能中微子望远镜（IceCube）因此可行。
\end{结论}

\begin{结论}{共振荧光与 Mössbauer：反冲能}
发射/吸收 $\gamma$ 光子时自由原子核反冲 $E_R=\frac{E_\gamma^2}{2Mc^2}$：$^{57}$Fe（$14.4$ keV）$E_R\approx2$ meV，远大于线宽 $5\times10^{-9}$ eV，不能共振吸收。
晶格中无反冲分数（Lamb--Mössbauer 因子）\重点{f=\ee^{-k^2\langle x^2\rangle}}，整块晶体承担反冲，可共振；可分辨 $10^{-13}$ 相对能量变化（测引力红移、超精细场）。原子光谱中反冲 $\frac{(\hbar k)^2}{2M}\sim$ kHz，同理产生激光冷却的反冲极限。
\end{结论}

\参考{Foot, \emph{Atomic Physics}；Herzberg, \emph{Molecular Spectra}；Krane, \emph{Introductory Nuclear Physics}；Griffiths, \emph{Introduction to Elementary Particles}；Kittel, \emph{Introduction to Solid State Physics}；PDG 综述；MIT 8.421、8.701、8.231。}
